\documentclass[../main.tex]{subfiles}
\begin{document}

\textit{This chapter describes the architecture, the data model, the API, and the design of the
three cross-cutting concerns that shaped the system most: authorisation, the AI cascade and the
offline policy.}

\section{Overall architecture}
Koshi is a client--server system. The browser runs a Vue single-page application that keeps a
local copy of the open project; an nginx container serves the application and forwards
\code{/api} to a FastAPI service, which owns the MySQL database and calls the AI provider. For
public access, a Cloudflare Tunnel connects outwards from the host to Cloudflare's edge, so no port
is opened on the router and TLS is terminated by Cloudflare (Figure~\ref{fig:arch}).

\begin{figure}[H]
\centering
\begin{tikzpicture}[
  box/.style={draw, rounded corners, align=center, font=\footnotesize, minimum height=0.9cm, fill=white},
  grp/.style={draw, dashed, rounded corners, inner sep=6pt},
  >=Stealth, node distance=0.5cm]
\node[box, minimum width=3.2cm] (ui) at (0,0) {Vue 3 components\\(board, table, Scrum, AI)};
\node[box, minimum width=3.2cm, below=of ui] (store) {Pinia task store\\\code{canMutate} gate};
\node[box, minimum width=3.2cm, below=of store] (idb) {IndexedDB\\(local copy)};
\node[grp, fit=(ui)(store)(idb), label={[font=\footnotesize\bfseries]above:Browser}] (br) {};

\node[box, minimum width=2.6cm] (cf) at (5,0) {Cloudflare edge\\TLS, DNS};
\node[box, minimum width=2.6cm, below=0.9cm of cf] (tun) {cloudflared\\tunnel};
\node[box, minimum width=2.6cm, below=0.9cm of tun] (ng) {nginx\\static + /api proxy};

\node[box, minimum width=3.2cm] (api) at (10,-0.2) {FastAPI routers\\auth, projects, tasks,\\sprints, stats, AI};
\node[box, minimum width=3.2cm, below=of api] (svc) {Services: AI cascade,\\notify, uploads, mentions};
\node[box, minimum width=3.2cm, below=of svc] (db) {MySQL 8.4\\(Alembic migrations)};
\node[grp, fit=(tun)(ng)(api)(svc)(db), label={[font=\footnotesize\bfseries]below:Docker host}] {};

\node[box, minimum width=2.6cm] (llm) at (10,2.0) {Groq API\\gpt-oss-120b};

\draw[<->] (ui) -- (store); \draw[<->] (store) -- node[right,font=\scriptsize]{write first} (idb);
\draw[->] (br.east) -- node[above,font=\scriptsize]{HTTPS} (cf);
\draw[<-] (cf) -- node[right,font=\scriptsize]{outbound} (tun);
\draw[->] (tun) -- (ng);
\draw[->] (ng.east) -- node[midway, sloped, above, font=\scriptsize]{REST/JSON} (api.west);
\draw[<->] (api) -- (svc); \draw[<->] (svc) -- (db);
\draw[->] (api.north) -- node[right,font=\scriptsize]{tier 1} (llm);
\end{tikzpicture}
\caption{Deployment and component architecture}\label{fig:arch}
\end{figure}

The frontend is organised by responsibility: \code{components/} (26 Vue components),
\code{stores/} (task, theme and language state), \code{lib/} (pure, framework-free modules:
topological sort and critical path, keyboard dispatcher, urgency ordering, git-diff parser, mention
parsing, translations), and \code{services/api.ts}, the only place that talks to the server. The
backend follows the same separation: \code{routers/} (HTTP), \code{schemas/} (Pydantic
contracts), \code{models/} (SQLAlchemy entities), \code{services/} (AI, notifications, uploads) and
\code{security.py} (authentication and the project-scope guards).

\section{Data model}
\begin{figure}[H]
\centering
\begin{tikzpicture}[ent/.style={draw, rounded corners, font=\scriptsize, align=left, fill=gray!6, minimum width=3.1cm}, >=Stealth]
\node[ent] (u) at (0,0) {\textbf{users}\\id, email, full\_name,\\hashed\_password, skills, avatar};
\node[ent] (pm) at (4.4,0) {\textbf{project\_members}\\project\_id, user\_id,\\role (PM/MEMBER),\\status (PENDING/ACCEPTED)};
\node[ent] (p) at (8.8,0) {\textbf{projects}\\id, name, description,\\owner\_id};
\node[ent] (s) at (8.8,-3) {\textbf{sprints}\\id, project\_id, name, goal,\\start/end\_date, is\_active,\\completed\_at, committed\_points};
\node[ent] (t) at (4.4,-3) {\textbf{tasks}\\id, project\_id, sprint\_id,\\status, priority, points,\\due\_date, completed\_at,\\dependencies\_json};
\node[ent] (ta) at (0,-3) {\textbf{task\_assignees}\\task\_id, user\_id};
\node[ent] (c) at (4.4,-6) {\textbf{comments}\\task\_id, author\_id, kind,\\content, parent\_id};
\node[ent] (a) at (8.8,-6) {\textbf{attachments}\\comment\_id, stored\_name,\\content\_type, size};
\node[ent] (e) at (0,-6) {\textbf{task\_events}\\task\_id, actor\_id, field,\\old\_value, new\_value};
\node[ent] (n) at (0,-8.4) {\textbf{notifications}\\user\_id, kind, actor\_id,\\task\_id, comment\_id, read\_at};
\draw (u) -- (pm); \draw (pm) -- (p); \draw (p) -- (s); \draw (p) -- (t); \draw (s) -- (t);
\draw (t) -- (ta); \draw (ta) -- (u); \draw (t) -- (c); \draw (c) -- (a); \draw (t) -- (e);
\draw (n.west) -- ++(-0.5,0) |- (u.west);
\end{tikzpicture}
\caption{Entities (10 tables)}\label{fig:er}
\end{figure}

Three design decisions matter most:
\begin{itemize}
  \item \textbf{Authority lives on the membership, not on the user.} \code{users} has no role
        column; \code{project\_members} holds the role and the invitation status. A PENDING
        membership grants nothing: the lookup used by every guard returns it only once accepted.
  \item \textbf{Several assignees per task} through \code{task\_assignees}; each must be an
        accepted member of the task's project.
  \item \textbf{History and Scrum need time.} \code{tasks.completed\_at} is set when a task
        enters DONE and cleared when it leaves, which is what the burndown reads;
        the sprint's \code{committed\_points} snapshots the commitment when a sprint starts, which is
        what velocity compares against. \code{task\_events} stores display strings rather than
        foreign keys, so an entry still reads correctly after the person or sprint it names is
        renamed or deleted.
\end{itemize}
The schema is created and changed only by nine Alembic migrations; outside development the
service refuses to start if the database is not at the latest revision.

\section{API design}
The REST API has 51 endpoints under \code{/api}; Table~\ref{tab:api} summarises them.

\begin{table}[H]
\centering
\caption{API overview}\label{tab:api}
\renewcommand{\arraystretch}{1.2}
\begin{tabular}{|p{2.9cm}|c|p{9.0cm}|}
\hline
\textbf{Router} & \textbf{Endpoints} & \textbf{Main operations} \\ \hline
\code{/auth} & 4 & register, login, Google sign-in, current user \\ \hline
\code{/users} & 5 & list with workload, update own profile, avatar upload/serve/delete \\ \hline
\code{/projects} & 11 & my projects, create, delete, members, invite, accept/decline, change role, remove \\ \hline
\code{/tasks} & 11 & list/filter, create, update, delete, cycle status, comments, attachments, history \\ \hline
\code{/sprints} & 9 & list, create, update, start, complete, delete, statistics, burndown, velocity \\ \hline
\code{/stats} & 2 & team workload, overdue tasks \\ \hline
\code{/notifications} & 4 & list, unread count, mark one / all read \\ \hline
\code{/ai} & 5 & weekly summary, meeting minutes, assignment, decompose, sprint retrospective \\ \hline
\end{tabular}
\end{table}

Errors follow HTTP semantics: 400 for an invalid reference (an unknown dependency, a sprint of
another project), 403 when a member attempts a PM action, 404 for a project the caller cannot
see---\emph{not} 403, because revealing that a project exists is itself a leak---and 409 when a
second sprint is started while one is active.

\section{Authorisation design}
Every project-scoped route calls one of two guards, \code{require\_member} or
\code{require\_project\_pm}, before touching data. The user interface hides controls a member cannot
use, but this is only an affordance: the server is the security boundary~\cite{owasp2021top10}.
During the project a code change proposed for the main branch made the client \emph{pretend} to be
PM after the server had refused a role change; it was rejected for exactly this reason. Sessions are
JWTs signed with a per-deployment secret that is generated on the host and never committed or
copied into an image; the startup check refuses weak or default secrets.

\section{AI cascade}
Each AI feature calls one function that tries three tiers in order (Figure~\ref{fig:cascade}).

\begin{figure}[H]
\centering
\begin{tikzpicture}[b/.style={draw, rounded corners, align=center, font=\footnotesize, text width=2.1cm, minimum height=1.1cm}, >=Stealth]
\node[b] (req) at (0,0) {AI request\\+ UI language};
\node[b] (t1) at (4.5,0) {Tier 1\\Groq (cloud)};
\node[b] (t2) at (9.0,0) {Tier 2\\Ollama (local)};
\node[b] (t3) at (13.5,0) {Tier 3\\deterministic};
\node[b] (val) at (6.75,-2.6) {Parse JSON, validate,\\normalise};
\draw[->] (req) -- (t1);
\draw[->] (t1) -- node[above=2pt, align=center, font=\scriptsize]{error /\\ timeout} (t2);
\draw[->] (t2) -- node[above=2pt, align=center, font=\scriptsize]{not\\ reachable} (t3);
\draw[->] (t1) -- node[left,font=\scriptsize]{200} (val);
\draw[->] (t2) -- (val);
\draw[->] (val) -- node[below right=1pt, font=\scriptsize, pos=0.45]{unusable} (t3);
\end{tikzpicture}
\caption{Three-tier AI cascade}\label{fig:cascade}
\end{figure}

The deterministic tier builds an answer from the project's own data (for example, the
retrospective counts committed and completed points, lists blocked tasks with their reasons and
detects scope added after the sprint started), so the feature still works with no model at all.
Falling through to tier 3 is logged as a degradation; this matters because the tiers return the same
shape and an expired key would otherwise look perfectly healthy. The interface language is sent
with each request (\code{Accept-Language}) and added to the prompt, so an English interface
receives English answers even when the tasks are written in Vietnamese.

\section{Offline policy and synchronisation}
The store applies every change to the screen and to IndexedDB first, then sends it to the server.
A single property, \code{canMutate}, decides whether a write is allowed: always when online; when
offline, only in a personal project. Every control that writes---including the AI tools that insert
tasks---checks it itself, because the store refuses silently and a control that did not check would
report success having written nothing. Two such defects were found and fixed during testing.

\section{User interface design}
The interface has three screens without a router: the landing page for visitors, the board, and
the profile page. Modals provide the task inspector, the Scrum panel and the AI tools. All visible
text comes from a translation table with an English and a Vietnamese entry for every key (more than
500 keys); a missing Vietnamese key is a compile error, and an automated test fails if a component
contains literal text.

\end{document}
